\newpage
\renewcommand{\sectioncolor}{chap4}
\section{Warstwa prezentacji – Implementacja funkcjonalności Frontend (React)}

Warstwa kliencka systemu została zbudowana przy użyciu biblioteki \textbf{React 19}. Poniżej opisano technologiczną realizację kluczowych modułów oraz funkcjonalności interfejsu użytkownika, skupiając się na logice biznesowej i integracji z mikroserwisami.

\subsection{Moduł Użytkownika – Proces rejestracji i rezerwacji}
Interfejs klienta został zaprojektowany z myślą o prostocie obsługi. Proces rejestracji obejmuje weryfikację adresu e-mail za pomocą cyfrowego kodu OTP, co zapewnia bezpieczeństwo tożsamości. Najważniejszym elementem jest formularz rezerwacji, który dynamicznie oblicza dostępność slotów czasowych na podstawie danych z serwisu rezerwacji.

\begin{lstlisting}[language=JavaScript, caption={Walidacja danych kontaktowych oraz filtrowanie slotów (BookingPage.js)}]
const validateContact = (next) => {
  const email = String(next.customerEmail || "").trim();
  const phone = sanitizePhoneNumberInput(next.customerPhone);

  const emailOk = !email || /^[^\s@]+@[^\s@]+\.[^\s@]+$/.test(email);
  const phoneOk = !phone || /^\+\d{8,15}$/.test(phone);

  return {
    customerEmail: emailOk ? "" : "booking.invalidEmail",
    customerPhone: phoneOk ? "" : "booking.invalidPhone"
  };
};

const filteredAvailableSlots = useMemo(() => {
  if (!form.date || !form.staffId) return [];
  return availableSlots.filter((s) => 
    String(s.staffId || "") === String(form.staffId)
  );
}, [availableSlots, form.date, form.staffId]);
\end{lstlisting}

\subsection{Zarządzanie profilem i Dashboard klienta}
Użytkownik po zalogowaniu ma dostęp do pulpitu (\textit{Dashboard}), który agreguje informacje o nadchodzących wizytach. System umożliwia również zarządzanie danymi profilowymi, wykorzystując zaawansowane funkcje normalizacji danych wejściowych, co gwarantuje spójność bazy danych.

\begin{lstlisting}[language=JavaScript, caption={Normalizacja kodu pocztowego i formatowanie numeru telefonu (AccountPage.js)}]
const normalizePostalCodeInput = (value) => {
  const digits = String(value || '').replace(/\D/g, '').slice(0, 5);
  if (digits.length <= 2) return digits;
  return `${digits.slice(0, 2)}-${digits.slice(2)}`;
};

const formatPhoneDisplay = (value) => {
  const raw = String(value || '').trim();
  const digits = raw.replace(/\D/g, '').slice(0, 12);
  if (!digits) return raw.startsWith('+') ? '+' : '';
  
  const country = digits.slice(0, 3);
  const rest = digits.slice(3);
  const groups = rest.match(/.{1,3}/g) || [];
  return `+${country} ${groups.join('-')}`;
};
\end{lstlisting}

\subsection{Moduł Firmowy – Zarządzanie uprawnieniami i kadrą}
Panel administracyjny firmy dynamicznie dostosowuje zakres dostępnych funkcjonalności na podstawie roli użytkownika (RBAC – \textit{Role-Based Access Control}). Logika frontendu rozróżnia uprawnienia właściciela (\textit{Owner}), managera oraz pracownika.

\begin{lstlisting}[language=JavaScript, caption={Logika weryfikacji uprawnień wewnątrz organizacji (CompanyPanel.js)}]
const companyRole = tokenManager.getUser()?.companyRole;
const userRoles = tokenManager.getUser()?.roles || [];
const isAdmin = userRoles.includes("Admin");
const isOwner = companyRole === "Owner";
const isManager = companyRole === "Manager";

// Definicja dostępu do akcji krytycznych
const canManageEmployees = isAdmin || isOwner || isManager;
const canEditCompany = isAdmin || isOwner;

const handleUpdateUserRole = async (userId, newRole) => {
  try {
    await companyUsersAPI.updateUserRole(userId, newRole);
    await loadEmployees();
  } catch (err) {
    console.error("Failed to update role", err);
  }
};
\end{lstlisting}

\subsection{Zarządzanie harmonogramami i przerwami pracowników}
Ważną funkcjonalnością jest wizualizacja czasu pracy. Użytkownik może definiować dni i godziny pracy oraz konfigurować przerwy (\textit{Staff Breaks}). Do obsługi wielodniowych harmonogramów wykorzystano komponent zarządzający kolekcją dni.

\begin{lstlisting}[language=JavaScript, caption={Obsługa komponentu wyboru dni tygodnia (CompanyPanel.js)}]
const toggleDay = (day) => {
  const d = Number(day);
  const current = normalizeDays(normalizedValue);
  const next = current.includes(d) 
    ? current.filter((x) => x !== d) 
    : [...current, d];
  onChange(normalizeDays(next));
};

const handleCreateStaffBreak = async (e) => {
  e.preventDefault();
  const daysToCreate = normalizeDays(staffBreakFilters.dayOfWeek);
  await Promise.all(
    daysToCreate.map((d) => staffBreaksAPI.create({ ...payload, dayOfWeek: d }))
  );
};
\end{lstlisting}

\subsection{Moduł Administracyjny i szablony powiadomień}
Administratorzy systemowi zarządzają szablonami wiadomości, które wspierają dynamiczne wstrzykiwanie zmiennych. Pozwala to na personalizację powiadomień e-mail i SMS wysyłanych przez mikroserwis powiadomień.

\begin{lstlisting}[language=JavaScript, caption={Mechanizm wstawiania tokenów do edytora szablonów (AdminPanel.js)}]
const insertTemplateToken = (token) => {
  const field = activeTemplateField === "subject" ? "subject" : "body";
  const el = field === "subject" ? subjectRef.current : bodyRef.current;
  const currentValue = String(editingTemplate?.[field] ?? "");

  if (el && typeof el.selectionStart === "number") {
    const start = el.selectionStart;
    const end = el.selectionEnd;
    const next = `${currentValue.slice(0, start)}${token}${currentValue.slice(end)}`;
    handleTemplateFormChange(field, next);
  }
};
\end{lstlisting}

\subsection{Integracja z geokodowaniem i mapami}
W widoku szczegółów usługi zaimplementowano interaktywną mapę \textbf{OpenStreetMap}. System przelicza adres tekstowy na współrzędne geograficzne i wyświetla lokalizację salonu w komponencie \textit{iframe}, co znacząco poprawia użyteczność aplikacji.

\begin{lstlisting}[language=JavaScript, caption={Geokodowanie adresu i generowanie widoku mapy (ServiceDetailsPage.js)}]
useEffect(() => {
  const geocode = async () => {
    try {
      const res = await geocodeAPI.geocode(osmQuery);
      if (res.data?.found) {
        setMapCoords({ lat: res.data.lat, lon: res.data.lon });
      }
    } catch (err) {
      setMapCoords(null);
    }
  };
  if (osmQuery) geocode();
}, [osmQuery]);
\end{lstlisting}